\section{Контракт гейтов и повторная проверка измерений}
\label{app:gates}

Численные результаты зафиксированного среза приведены в
разделе~\ref{sec:olmo}. Ниже задаются условия, при которых измерение
вообще может обосновать соответствующее утверждение. Единица сравнения ---
один чекпойнт, один токенизатор и один фиксированный интерфейс.
Точное описание редакции хранится в
\src{observer/gate_contract_ru.md}{контракте второй серии}.

\paragraph{Панель и бюджет.}
Восемь битов входа записываются как $\mathbf{x}=LwR$.
Середины $w\in\{0,1\}^6$ сортируются по возрастанию SHA-256 от ASCII-строки
\verb!observer-v2|! с приписанными шестью битами, включая начальные нули.
При нумерации с единицы ранги 1--8 отведены для dev, 9--20 --- для
совместной проверки G1 и G2, 21--32 --- для G3, 33--44 --- для G4,
45--64 образуют резерв. Для каждой середины предъявляется весь квадрат
концов и весь квадрат историй: $(L,R,h,e)\in\{0,1\}^4$ в лексикографическом
порядке. Это 16 строк на блок, 128 строк на dev, 192 на каждую указанную
12-блочную панель. На dev подбирают формулировку, демонстрации, слой и
масштаб вмешательства. Демонстрационные строки отделены от проверочных.
Сравниваемые реализации замораживают до открытия их общей отложенной панели.
Для метода, выбранного по результатам открытой панели, требуется свежая
панель.

\paragraph{Получение обеих историй.}
Сначала измеряются полнословарные вероятности однотокеновых продолжений
\texttt{[пробел]A}, \texttt{[пробел]B} при температуре 1.
Для допуска нужны $P(A),P(B)\ge 0{,}01$. Выполняются ровно 1024 выборки
с возвращением из условного бинарного распределения
$P(A\mid A\text{ или }B)=P(A)/(P(A)+P(B))$, генератором
\texttt{random.Random(110)}. Сохраняются обе исходные вероятности, вся
последовательность выборок и позиции первых A и B. В каждой истории модель
обрабатывает первый фактически выбранный токен соответствующего типа.
Снимок снимается после обработки этого токена и до обратной связи. Если
минимум вероятности или присутствие обеих букв не обеспечены, для этого
префикса выносится \texttt{NOT\_EVALUABLE}. Принудительная подстановка
буквы и дополнительная серия сэмплирования изменили бы условие опыта.

\paragraph{Общие критерии.}
На каждой обязательной числовой строке, в том числе на нейтральных,
восстановленных и подменённых контролях, сохраняются цель, полный
словарный максимум, $\ell_1-\ell_0$ и $q=P(0)+P(1)$. Вероятности
вычисляются после полнословарной нормировки в FP32 или выше. Требуются
конечные значения и $q\ge 0{,}80$ на каждой строке. При строгом знаке
зазора победивший бит имеет вероятность $>0{,}40$, а любой посторонний
токен --- не более $0{,}20$. При нулевом зазоре строка остаётся в
знаменателе и засчитывается как неуспех. Полный блок требует 16/16
правильных выходов. Рабочие пороги составляют 6/8 полных блоков на dev
и 9/12 на отложенной панели. Отдельно публикуются парные решения при
$L\ne R$ для каждого $e$. Указанные 75\% относятся к доле целых блоков,
а повторное использование одного префикса не образует независимую выборку
экземпляров модели.

\paragraph{G0 и G0R.}
G0 сверяет хэши весов и токенизатора, однотокеновость A/B/X и
0/1 непосредственно в их контексте, точную композицию идентификаторов,
форму, тип данных и позиции RoPE. Диагностики работают на одноразовых
клонах. Чтение прошлого выбора останавливается до обратной записи ответа.
При наличии и при отсутствии этой диагностики один и тот же основной запрос
должен дать одинаковые полные логиты. Исходные снимки обеих ветвей сохраняют
байты всех активных K/V, смещение и порядок видимых позиций, включая индексы
кольцевого буфера. Сравниваются логические данные активных позиций. Свободная
ёмкость аллокатора вне этого диапазона в проверку не входит.

G0R независимо от поведенческого счёта выполняет до отбора обратимые подмены
позиции выбора во всех слоях. При общей предыстории $U$ снимки имеют вид
$C_A=(U,a)$, $C_B=(U,b)$. Для каждого текста сигнала и каждой новой строки
тестируются маршрут $A\to B\to A$ и зеркальный маршрут из B.
Обратная подмена применяется к результату прямой; считывания выполняются на
его клонах. В средней точке сравнивается натуральный донор, в конечной ---
исходная история. Сверяются метаданные, активные K/V и логиты при одинаковой
раскладке. Формы, конечность и байты исходного типа должны совпасть, а после
расширения логитов до FP32 требуется максимум абсолютной разности,
равный нулю. Эталонные снимки сверяются до и после. Для первого Base-прогона
зафиксирован исходный \texttt{prefix(8)}. Перед подменой проверяются
идентичность общего префикса и существование хотя бы одного различающегося
активного байта позиции выбора. Если ветви неразличимы уже в этой позиции,
результат G0R для данного префикса --- \texttt{NOT\_EVALUABLE}.
На dev при 128 строках два сравнения логитов в каждой строке дают 256
проверок.

\paragraph{G1 и G2.}
G1 проверяет текстовое вычисление на натуральных историях. Помимо
обязательных числовых проб на независимых копиях измеряются D1 (воспроизвести
$h$ до обратной связи), D2 (назвать исправленное правило для четырёх пар
$h,e$ без новой строки) и D3 (применить явно названное правило в свежем
эпизоде). Буквенные диагнозы публикуются отдельно от основной числовой
задачи. Успех текстового контроля определяется форматом и долей целых
блоков. Для G2 после G0, G0R и G1 на панели 9--20 повторяются сравнения
донор/возврат; добавляется общий кэш с нейтральным токеном X в той же
позиции и той же длины. Равенство с натуральным донором следует из
$P_B(U,a)=(U,b)$. Поэтому положительный технический вердикт есть
\texttt{REPLAY\_OK}. Поведенческий эффект по натуральным историям
уже оценён G1.

\paragraph{G3.}
После G1 и \texttt{REPLAY\_OK} текстовый сигнал заменяется одним
вмешательством до записи K/V на нейтральном токене. До теста задаются токен,
позиция, слой, точка вмешательства, источник вектора и масштаб. В типе
операции представляется один конечный ненулевой $v$. Импульсы $v_0=v$, $v_1=-v$
получаются точной сменой знака. Для одного $e$ обе истории получают
побайтно один тензор. Проверяются форма, тип, поэлементная противоположность
и равенство ненулевых евклидовых норм при одинаковом порядке вычисления в
FP64. Отдельно для каждой ветви записывается норма фактически округлённого
приращения $\operatorname{round}(r+v_e)-r$. После одного сложения
вмешательство выключается. При нулевом вмешательстве обе ветви $e$ с
фиксированными $h,\mathbf{x}$ должны иметь идентичные входы и логиты;
публикуются их точность и масса. Для положительного результата с импульсом
требуются общий формат и минимум 9/12 полных блоков на панели G3.

\paragraph{G4.}
До предъявления $\mathbf{x}$ четыре донора обрабатывают общий нейтральный
промежуток после $h,e$. В свежий шаблон переносятся K/V исключительно
заранее заданных позиций и слоёв этого промежутка, обозначаемые $M_{he}$.
Все остальные активные позиции, токены, длины и метаданные идентичны у
четырёх получателей. Прямой доступ к $h,e$ и иным потомкам донорской
истории закрыт. Нейтральное состояние $M_*$ единожды определяется общей
нейтральной историей. Для текстового канала требуются G1 и G2, для
импульсного также G3. Условия на 12 серединах совместны:
\begin{enumerate}
\item Получатель с полной историей и тем же промежутком проходит формат и
решает минимум 9/12 полных блоков. При неуспехе этого контроля отрицательный
перенос не устанавливает отсутствия нужной информации в доноре.
\item После установки $M_*$ все четыре входа, активные K/V и логиты при
фиксированном $\mathbf{x}$ побайтно совпадают. На одном производном снимке
возврат своего $M_{he}$ после чужого восстанавливает исходные логиты.
\item При переносе своего $M_{he}$ выполнены формат и порог 9/12 полных
блоков. Каждый блок включает всех четырёх доноров и все четыре сочетания
концов.
\item На 96 строках с $L\ne R$ обозначим число верных решений при переносе
и при стирании через $C_M$ и $C_*$. Требуются валидный формат,
$C_*\le48$, 0/12 полных блоков после стирания и
$(C_M-C_*)/96\ge 0{,}25$. Исправленные и испорченные строки учитываются
раздельно. При обязательном одном строгом бинарном ответе общий вход даёт
ровно 48/96. Девять полных блоков при переносе обеспечивают по меньшей
мере 72/96 на расходящихся концах.
\end{enumerate}

\paragraph{G4-1.}
Отправитель до предъявления какой-либо $\mathbf{x}$ свободно выбирает один
токен полнословарным greedy. Допустимы только два заранее заданных
однотокеновых окончания A/B. Ничья за первое место или другой токен даёт
\texttt{FAIL\_FORMAT} без повторной попытки. Четыре сообщения для историй
$00,01,10,11$ фиксируются в этом порядке. В исходной спецификации для Base
зафиксированы \texttt{prefix(8)}, буквальный сигнал correct/wrong и запрос
на запись исправленного правила. У измеренного двухшагового Instruct
собственные промпты и дополнительный порог $P(A)+P(B)\ge0{,}80$
зафиксированы в \src{observer/two-step-interface-04.json}{манифесте варианта A}.
Условия двух чекпойнтов учитываются отдельно.

Для воспроизведения именно исходной спецификации Base отправителю после
выбора передают \verb!\nFeedback: correct! либо \verb!\nFeedback: wrong!,
затем \verb!\nRecord the corrected rule. Rule:!. В этих записях
\verb!\n! обозначает один перевод строки. Свежий префикс получателя
состоит из двух строк:
\begin{verbatim}
Rule A selects the first bit. Rule B selects the last bit.
Rule:
\end{verbatim}
Сразу после двоеточия вставляют выбранный ID с начальным пробелом. За ним
идут \verb!\nBits: !, восемь битов через одиночные пробелы и
\verb!\nAnswer: ! с пробелом в конце.

В свежем эпизоде получателю передаётся только выбранный ID, для которого
он заново вычисляет K/V. Токены общего префикса и строки запроса сверяются
с токенизацией целого текста. Во внешний интерфейс не входят K/V или логиты
отправителя, номер ветви, время ответа и параметры сэмплирования. Выбор
носителя или ветвей по $h,e,p$, правильным ответам либо их заменителям
запрещён. Для стирания подаются те же входы с одним заранее фиксированным
токеном A во всех мирах. Для \texttt{ZERO\_ERROR\_ON\_PANEL} требуются
четыре допустимых сообщения, проверки границы канала и стирания,
числовой формат и 12/12 полных блоков. При совместном испытании G4 и G4-1
обе версии замораживаются до первого открытия панели 33--44.

\paragraph{Приоритет статусов и воспроизводимость.}
Техническая ошибка, нечисловое значение или нарушение требуемой побайтовой
идентичности дают \texttt{INVALID}. Затем по убыванию приоритета следуют
\texttt{NOT\_EVALUABLE}, \texttt{FAIL\_FORMAT} и
\texttt{FAIL\_DEV}/\texttt{FAIL\_TEST}. Зависимая проверка, до которой не
дошли, получает \texttt{NOT\_RUN}. Новая редакция интерфейса или другой
чекпойнт образуют отдельное условие.

Из корня зафиксированного комплекта команды
\begin{verbatim}
python3 -B observer/verify.py
python3 -B observer/verify_part_iii.py
\end{verbatim}
решают разные задачи. Первый валидатор пересчитывает конечные конструкции,
проверяет манифест исторического корпуса в публикационной редакции,
адреса реестра и восстановление таблицы из SVG. Второй читает сохранённые
dev-отчёты, проверяет их сетки условий, записанные вероятности и логиты,
получает точность, целые блоки, противоположные пары, диагностики и эффект
стирания. Для G0/G0R он повторно сверяет метаданные и зарегистрированные
хэши попарных сравнений нативных логитов. Перезапуск на GPU и повторное
побайтное сравнение заново вычисленных тензоров составили бы новое
измерение.
